% Generated draft: requires mathematical and source review.
\documentclass[11pt]{article}
\usepackage[margin=25mm]{geometry}
\usepackage{fontspec}
\setmainfont{Latin Modern Roman}
\usepackage{amsmath,amssymb,mathtools,unicode-math}
\setmathfont{Latin Modern Math}
\usepackage{enumitem,needspace,xurl,hyperref}
\setlength{\parindent}{0pt}
\setlength{\parskip}{5pt}
\setlength{\emergencystretch}{3em}
\newcommand{\N}{\mathbb{N}}
\newcommand{\Z}{\mathbb{Z}}
\newcommand{\Q}{\mathbb{Q}}
\newcommand{\R}{\mathbb{R}}
\newcommand{\C}{\mathbb{C}}
\newcommand{\F}{\mathbb{F}}
\newcommand{\A}{\mathbb{A}}
\newcommand{\PP}{\mathbb{P}}
\newcommand{\E}{\mathbb{E}}
\newcommand{\eps}{\varepsilon}
\newcommand{\dd}{\,\mathrm{d}}
\begin{document}
\section*{UnsolvedMath ID: 700726}
\textbf{Source problem code:} catalogue entry
\subsection{Definitions and mathematical statement}
Let \(q\) be a prime power and let \(G(q)\) be a simple Chevalley group over \(\F_q\). For a positive integer \(m\), write \(G(q^m)\) for the corresponding group of the same type over \(\F_{q^m}\). The source presupposes an inclusion \(\iota_m:G(q)\hookrightarrow G(q^m)\) arising from the field extension; its compatibility with the chosen simple-group quotients requires clarification. For a field \(k\), a representation means a \(k\)-vector space with a linear group action. An irreducible representation is a nonzero representation with no nonzero proper invariant subspace. For a module of finite length, its composition factors are the simple quotients in a composition series. Write \(\operatorname{Irr}_k(H)\) for the isomorphism classes of irreducible \(k\)-representations of \(H\).

The supplied problem, part (b), asks whether there exists a positive integer \(m\) such that, for every field \(k\) of prime characteristic \(\ell\) with \(\ell\nmid q\), the restriction to \(G(q)\) of every non-one-dimensional representation of \(G(q^m)\) has every member of \(\operatorname{Irr}_k(G(q))\) as a composition factor.

For finite-dimensional representations, the literal reading of this assertion is
\[
\exists m\in\Z_{\geq1}\ \forall k\ \forall V:\quad
\left[\operatorname{char}(k)=\ell>0,\ \ell\nmid q,\ V\text{ is a finite-dimensional }kG(q^m)\text{-module},\ \dim_k V\ne1\right]
\Longrightarrow
\left[\text{every }W\in\operatorname{Irr}_k(G(q))\text{ occurs in }\operatorname{Res}_{\iota_m}V\right].
\]
Here occurrence means occurrence as a composition factor, without a prescribed multiplicity. This display is a conditional formalization under the finite-dimensional convention, which the source does not explicitly state. The source also does not explicitly require \(V\) to be irreducible, to have nontrivial action, or to be nonzero. Consequently, its phrase ``non-one-dimensional'' leaves material formulation issues unresolved; no irreducibility hypothesis has been silently inserted. The displayed quantifier order takes \(m\) to be independent of \(k\) and \(\ell\); this scope should be checked against the original source.

The supplied archived part (a) asks for the analogous assertion over \(\C\): there exists a positive integer \(m\) such that the restriction of every non-one-dimensional complex representation of \(G(q^m)\) to \(G(q)\) contains every irreducible complex representation of \(G(q)\) as a composition factor. The same unresolved qualification concerning the representation class applies. The supplied archive reports part (a) as proved; part (b) remains in the supplied main section. These are supplied source classifications, not independently verified mathematical or literature assessments.
\subsection{Short English statement}
For each simple Chevalley group \(G(q)\), part (b) asks for one extension degree \(m\) such that every representation described by the source as non-one-dimensional, in prime characteristic not dividing \(q\), contains every irreducible representation of \(G(q)\) after restriction. The complex analogue, part (a), is supplied as archive-solved. The intended representation class, coefficient-field conventions, and subgroup inclusion need clarification.
\subsection{Sources}
\begin{itemize}
\item Local UnsolvedMath registry, permanent ID \(700726\), exact source code \texttt{catalogue entry}, title ``Kourovka Notebook Problem 11.33''; supplied upstream revision \texttt{d371ee917ac9647fbb06cd6fc529deec6f3dda82}.
\item \emph{The Kourovka Notebook (Unsolved Problems in Group Theory)}, 11th issue (1990), Problem 11.33; supplied background identifies an update in the 21st issue (2026), \texttt{arXiv:1401.0300v48}. The supplied collection description names E. I. Khukhro and V. D. Mazurov as editors. These references have not been independently consulted.
\item Proposed in 1990 by the person recorded as \texttt{A. E. Zalesski\textbackslash u\{\textbackslash i}; the supplied name encoding is incomplete and has not been reconstructed.
\item The supplied Archive of Solved Problems entry attributes the proof associated with archived part (a) to D. Gluck, \emph{J. Algebra} \textbf{155}, no. 2 (1993), 221--237. No article title or external link is supplied, and the citation and its precise scope have not been independently verified.
\item Supplied status review dated \texttt{2026-10-16}: canonical status \texttt{partially\_solved}, evidence label \texttt{published}; the review says that solved parts occur in the archive and remaining parts in the main section. This is supplied dated curation, not an independently verified assessment.
\end{itemize}
\end{document}
